\documentclass[11pt,a4paper]{article}
\usepackage[utf8]{inputenc}
\usepackage[portuguese]{babel}
\usepackage{geometry}
\usepackage{fancyhdr}
\usepackage{listings}
\usepackage{xcolor}

\geometry{left=2cm,right=2cm,top=2cm,bottom=2cm}
\setlength{\parskip}{0.4em}
\setlength{\headheight}{14pt}
\pagestyle{fancy}
\fancyhf{}
\lhead{CC8550 --- Simulacao e Teste de Software}
\rhead{Atividade 09}
\cfoot{\thepage}

\lstset{
    basicstyle=\ttfamily\small,
    breaklines=true,
    frame=single,
    backgroundcolor=\color{gray!10}
}

\title{
    {\Large \textbf{Centro Universitario FEI}}\\
    \vspace{0.3cm}
    {\large Ciencia da Computacao}\\[0.8cm]
    {\LARGE \textbf{Simulacao e Teste de Software}}\\[0.3cm]
    {\Large Atividade 09 --- Sistema Task Manager}\\[0.8cm]
    {\large \textbf{Testes orientados a objetos e componentes}}
}
\author{
    \textbf{Aluno:} Lucas Reboucas Silva\\
    \textbf{R.A.:} 22.122.048-6\\[0.4cm]
    \textbf{Professor:} Luciano Rossi
}
\date{1o Semestre de 2026}

\begin{document}
\maketitle
\thispagestyle{empty}
\newpage
\setcounter{page}{1}

\section{Objetivo}
Implementar um sistema simples de gerenciamento de tarefas em Python e validar seu comportamento com testes automatizados de unidade e de componente, conforme os conceitos da Aula 09.

\section{Implementacao}
O projeto foi organizado em tres partes principais:
\begin{itemize}
    \item \texttt{Task}: modela uma tarefa com id, titulo, descricao, prioridade, prazo e status.
    \item \texttt{InMemoryStorage}: armazena tarefas em memoria com operacoes de adicionar, buscar, listar e remover.
    \item \texttt{TaskRepository}: camada de persistencia que atribui IDs e delega acesso ao storage.
\end{itemize}

Na classe \texttt{Task}, o metodo \texttt{validar()} garante:
\begin{itemize}
    \item titulo com pelo menos 3 caracteres;
    \item prazo nao pode estar no passado.
\end{itemize}
Tambem foi implementada verificacao para impedir atribuicao de status fora do enum \texttt{Status}.

\section{Tipos de testes}

\subsection{Teste de Unidade}
\textbf{Definicao}: Testa uma classe ou funcao isolada, sem dependencias externas. \\
\textbf{Foco}: Validar logica interna e regras de negocio. \\
\textbf{Caracteristica}: Rapido, nao requer banco de dados ou chamadas HTTP. \\
\textbf{Exemplo no LAB-09}: Testes da classe \texttt{Task} validam validacoes de titulo e prazo isoladamente.

\subsection{Teste de Componente}
\textbf{Definicao}: Testa um componente que colabora com dependencias externas, mas essas dependencias sao mockadas (simuladas). \\
\textbf{Foco}: Validar interacao entre componentes e delegacao correta de responsabilidades. \\
\textbf{Caracteristica}: Mais rapido que teste de integracao (pois usa mocks), mas mais lento que unitario. \\
\textbf{Exemplo no LAB-09}: Testes de \texttt{TaskRepository} mockam \texttt{storage} para validar se repository chama storage corretamente.

\section{Estrategia de testes}

\subsection{Testes unitarios (Task)}
Nos testes de \texttt{Task}, a classe e exercitada de forma isolada, sem mocks.

\subsubsection{Teste 1: Estado inicial da tarefa}
Instancia \texttt{Task} com parametros validos e verifica se atributos foram atribuidos corretamente e status inicial e \texttt{PENDENTE}.

\subsubsection{Teste 2: Validacao de titulo muito curto}
Tenta criar tarefa com titulo de 1-2 caracteres. Usa \texttt{pytest.raises(ValueError)} para confirmar que \texttt{validar()} lanca excecao com mensagem esperada. Testa regra: titulo minimo 3 caracteres.

\subsubsection{Teste 3: Validacao de prazo no passado}
Tenta criar tarefa com prazo anterior a \texttt{datetime.now()}. Valida que \texttt{validar()} rejeita prazos retroativos, testando restricao de logica de negocio.

\subsubsection{Teste 4: Mudanca de status valida}
Cria tarefa com status \texttt{PENDENTE} e altera para \texttt{EM\_PROGRESSO}. Verifica se property setter permite transicao e estado muda corretamente.

\subsubsection{Teste 5: Rejeicao de status invalido}
Tenta atribuir valor invalido (string ou numero) para status. Usa \texttt{pytest.raises} para confirmar que setter rejeita dados nao-enum com \texttt{ValueError}.

\subsection{Testes de componente (TaskRepository)}
Nos testes de \texttt{TaskRepository}, a logica interna do repositorio e real, enquanto a dependencia externa (\texttt{storage}) e mockada usando \texttt{unittest.mock.Mock}.

\subsubsection{Teste 6: Save atribui ID sequencial}
Chama \texttt{repo.save(task)} e valida se a tarefa retornada tem \texttt{id = 1}. Testa se repositorio gerencia auto-incremento.

\subsubsection{Teste 7: Save delega para storage.add()}
Apos \texttt{save()}, verifica se mock foi chamado exatamente uma vez com argumentos \texttt{(1, task)} usando \texttt{mock.add.assert\_called\_once\_with()}. Testa delegacao correta.

\subsubsection{Teste 8: Find by ID usa storage.get()}
Configura mock para retornar tarefa e chama \texttt{find\_by\_id(1)}. Valida se resultado e a tarefa mockada, testando consulta por ID.

\subsubsection{Teste 9: Sequencia save e find preserva integridade}
Salva tarefa, moqueia storage para retornar a mesma tarefa, depois busca. Valida que save e find trabalham juntos corretamente. Testa fluxo real de persistencia.

\section{Resultados}
Comando utilizado:
\begin{lstlisting}
pytest -v
\end{lstlisting}

Resultado esperado: todos os testes de \texttt{tests/test\_task.py} e \texttt{tests/test\_repository.py} passando, cobrindo os requisitos obrigatorios da atividade.

\section{Conclusao}
O laboratorio permitiu aplicar, na pratica, testes de estado, fixtures (setup), isolamento por mock/stub e validacao de colaboracao entre metodos. A separacao entre teste unitario e teste de componente ficou clara na diferenca entre \texttt{Task} (isolada) e \texttt{TaskRepository} (com dependencia externa mockada).

\end{document}
